%!TEX root=main.tex

\section{Shallow wPRFs in Higher-Level MPC Protocols}
\label{sec:applications}


The existence of pseudorandom functions computable by low-depth circuits has a long history, with strong ties to fundamental questions in complexity theory such as circuit lower bounds~\cite{STOC:RazRud94,C:MilVio12,SerVio12}, derandomization~\cite{FOCS:NisWig88,STOC:Williams13}, and learning theory~\cite{valiant1984theory,kearns1994cryptographic,C:BCGIKS21}. These connections are not the focus of this work and will not be covered, as they are typically of little relevance to cryptanalysis and symmetric cryptography.\footnote{We do note, however, that the tools and results developped in these works have significantly influenced the design of shallow wPRFs that we will cover here. Therefore, on several occasions, we will refer to some of these early works to explain the security rationale underlying some candidates.} In this work, we focus on the use of low-depth weak pseudorandom functions in practically-minded higher-level secure computation protocols (we will use the terms secure computation and MPC, short for MultiParty Computation, interchangeably). More precisely, we will be mostly concerned with cryptographic constructions that use a wPRF as an \emph{inner component} within an \emph{outer component}, where the latter is a high-end cryptographic building block (a primitive or a protocol) that uses internally the \emph{circuit description} of the wPRF (that is, we do not consider cryptographic protocols that make a ``black-box'' use of a wPRF, since such protocols are not sensitive to low-level details of the computational model where the wPRF is computed). In equivalent terms, these constructions compile a wPRF whose circuit description meets specific structural constraints into a cryptographic primitive that fulfills a set of security and efficiency requirements.


 Unlike the traditional study of shallow wPRFs, whose focus was generally on low-depth wPRFs in standard computational models of complexity theory\footnote{That is, the usual ``low-depth'' classes found in complexity-theory textbooks: $\AC^0$, $\AC^0[\oplus]$, $\ACC^0$, $\TC^0$, and related classes. They capture functions computable by constant-depth circuits operating with a certain set of gates, e.g., unbounded fan-in AND and OR gates for $\AC^0$.},  the computational models that best capture these target applications are slightly more exotic at first sight. Indeed, as will become clear next, there is a duality between the computational model in which a wPRF is implemented, and the class of functions that the outer cryptographic primitive or protocol can ``evaluate'' efficiently.

\subsection{Why Weak Pseudorandom Functions?}

Pseudorandom functions turn a short random key into a virtually unbounded number of pseudorandom strings. Advanced cryptographic protocols often consume a tremendous amount of randomness, as randomness is unavoidable to hide secret data used in complex interactions, but also to guarantee many other standard security properties (e.g., catching cheaters with random challenges, authenticating data with MACs...). As such, that pseudorandom functions can be leveraged in some of these higher-level protocols is not too surprising.

What is perhaps less natural is the focus on \emph{weak} pseudorandom functions: by definition, wPRFs require the participants to be given access to a trusted source of randomness in the first place, as pseudorandomness holds only when the inputs are truly random. However, it turns out that wherever PRFs are used as an inner component in a higher-level MPC protocol, wPRFs also suffice, for in all such applications, the inputs can be easily \emph{preprocessed} (either by the party holding them, or by all participants when they are public). When one is allowed to preprocess inputs, there is an easy way to convert any wPRF into a PRF: given an arbitrary input $x$, hash $x$ into $\Hash(x)$ using your favorite hash function (e.g., $\mathsf{SHA}$-256) before feeding it to the wPRF. This general methodology, dubbed the hash-then-evaluate paradigm in~\cite{SCN:BCEKM24}, belongs to the large family of heuristic hash-based transforms that are provably secure when the hash is modeled as a random oracle (similar to the Fiat-Shamir transform, the Fischlin transform, the Fujisaki-Okamoto transform, and many more). In addition, the work of~\cite{SCN:BCEKM24} showed that for several wPRF candidates (in fact, for essentially those covered in this SoK), instantiating $\Hash$ with a PRF whose key is publicly known yields a candidate PRF whose security guarantees are comparable to that of the original wPRF.

Therefore, if one is fine with relying on the random oracle methodology (for higher-level MPC protocols targeting real-world applications, this is usually the case) or on the security arguments outlined in~\cite{SCN:BCEKM24}, using a wPRF suffices. Furthermore, building shallow wPRFs turns out to be considerably easier than building shallow PRFs. This is a known general behavior both in theoretical cryptography\footnote{For example, PRFs cannot exist in $\AC^0$, even by using a constant number of xor and majority gates~\cite{SODA:Viola13}, while wPRFs in $\AC^0$ exist under standard assumptions such as factoring~\cite{STOC:Kharitonov93}.} and in concrete candidates design. In particular, all candidate wPRFs covered in this SoK are provably not PRFs (they are easily broken given chosen queries), and constructing shallow PRFs of a comparable complexity is an open question (or, in the case of Goldreich's PRG, is known to be impossible~\cite{SODA:Viola13,TCC:AppRay16b}).

\paragraph{On computational models and cost metrics.} There are two types of restrictions that can be enforced by the computational model, that we will losely refer to as ``hard'' and ``soft'' restrictions:
\begin{itemize}
	\item An outer scheme induces a computational model with a hard restriction when a measure of the circuits in the model (e.g., depth, gate type) cannot exceed a fixed bound. As a simple example, using a linearly-homomorphic encryption scheme (e.g., Paillier encryption~\cite{EC:Paillier99}) as outer component restricts the inner component (the primitive being homomorphically evaluated) to be computable by a circuit computing a linear combination over a ring; using instead the Boneh-Goh-Nissim encryption scheme~\cite{TCC:BonGohNis05} restricts the inner component to be computable by a \emph{depth-1} arithmetic circuit with product gates of indegree 2 (i.e., a quadratic polynomial).
	\item In other settings, the outer scheme induces a soft restriction on the circuits, where some measures (e.g., depth, gate count) can grow arbitrarily but are associated to a \emph{cost} (an efficiency metric of the final protocol, typically computation, communication, or round complexity) that increases as the measure increases. A typical example would be the use of BGV-style fully-homomorphic encryption~\cite{STOC:Gentry09,ITCS:BraGenVai12}: the inner component can be any Boolean circuit, but the computational cost (measured by noise growth, or number of necessary bootstrappings) grows with the multiplicative depth of the circuit.
\end{itemize}
Below, we cover the main target applications of these cryptographic constructions. For each target application, we provide some motivation, list the high-end cryptographic building blocks used as outer component to realize the application, and extract the computational model that captures the wPRFs that can be evaluated within each outer component, indicating whether it comes with soft or hard restrictions on the type of circuits that can be computed.

\subsection{Oblivious Pseudorandom Functions}
\label{subsec:oprf}

An oblivious pseudorandom function (OPRF)~\cite{TCC:FIPR05} for a pseudorandom function family $\PRF = (\KeyGen, \Eval)$ is a two-party protocol where one party holds a key $\msk$ in the support of $\KeyGen$, and the other party holds an input $x$ to the PRF. At the end of the protocol, the party holding $x$ should learn $\Eval(\msk,x)$ (and nothing more), while the party holding $\msk$ should not learn anything. OPRFs enjoy many applications in settings such as obliviously searching databases~\cite{TCC:FIPR05}, private set intersection and pattern matching~\cite{TCC:HazLin08,CCS:KKRT16}, password-protected secret sharing and password-authenticated key exchange~\cite{AC:JarKiaKra14}, single sign-on with privacy~\cite{EUROSP:BFHLY20}, cloud key management~\cite{CCS:JarKraRes19}, and many more---see~\cite{EUROSP:CasHesLeh22} for a recent survey on the topic.

\subsubsection{From General Secure Computation}

The simplest way to build an OPRF is to take a PRF family $\PRF = (\KeyGen, \Eval)$, and to run a secure 2-party protocol that, on input $\msk$ from one party and $x$ from the other, securely computes $\Eval(\msk,x)$. Secure $2$-party computation (2PC) refers to a general set of techniques involving two participants with private inputs $x$ and $y$ who wish to compute a public function $f(x,y)$ on their joint private inputs (with the output being revealed to one participant, or both). The two main paradigms for modern secure 2-party computation are secret-sharing-based 2PC in the preprocessing model~\cite{STOC:GolMicWig87,AC:FKOS15,CCS:KelOrsSch16,C:DPSZ12,EC:KelPasRot18} and garbled-circuit-based secure computation~\cite{FOCS:Yao86,ICALP:KolSch08,EC:ZahRosEva15,CCS:WanRanKat17b,CCS:WanRanKat17a,C:DILO22,EC:CWYY23}, the latter being typically more expensive but requiring much less rounds of communication than the former, which can make it a better option over high-latency networks. Secure computation is a very active topic: there are many available protocols, and the cost metric can vary significantly from one protocol to the other. The metrics are typically of the ``soft restriction'' type: 2PC can theoretically evaluate any function. Nevertheless, the cost is mostly driven by the following properties.
\begin{description}
\item[Free Linear Operations.] Linear operations (e.g., XORs) are typically ``for free'' in such systems, either because they can be performed locally by the participants in secret-sharing-based 2PC, or due to the free-XOR trick for garbled circuits~\cite{ICALP:KolSch08}. In contrast, the number of nonlinear gates (e.g., AND gates) matters.
\item[Low Depth.] In secret-sharing-based 2PC, the depth of the circuit (the length of the largest path from an input to an output, counted ignoring the ``free'' linear gates) is a major measure of efficiency, as it translates to a lower bound on the number of messages the participants must exchange. For example, when written as a Boolean circuit with XOR gates and AND gates, the \aes{} block cipher has depth 40~\cite{EC:ARSTZ15}. When the protocol takes place between distant machines, this incurs a significant latency: a single round trip between relatively close machines (e.g., in the same country) adds 50–100ms~\cite{SP:DKLS19}. For distant machines, latency can range from 100–260ms, up to 0.5 seconds in some cases~\cite{SP:DKLS19}. A 200ms roundtrip would result in 8s of delay for the \aes{} circuit.
\item[Correlated Randomness.] In the preprocessing model, the participants initially run a secure protocol to distribute correlated randomness among them. Different types of correlated randomness can be used to implement efficiently some ``complex'' gates. Typical examples include gates converting between arithmetic values over different fields~\cite{TCC:BIPSW18}, matrix operations~\cite{C:BCGIKS20}, or some standard non-linear functions used in machine learning such as equality tests, threshold predicates, and ReLU~\cite{TCC:BoyGilIsh19,EC:BCGGIK21}.
\end{description}
As the methods and protocols for efficiently distributing complex forms of correlated randomness evolve rapidly, it is not feasible to pinpoint a specific computational model and cost metric that would faithfully capture the evolving state-of-the-art. Nevertheless, a good rule of thumb is that secure computations shine when evaluating \emph{low-depth} circuits that minimize the number of non-linear gates, and can rely on certain types of more complex gates such as equality tests, greater-than predicates, ReLU, splines, or linear algebra, among others.

The performance gain obtained by using dedicated wPRFs to build OPRFs is significant. In Section~\ref{sec:bipsw}, we introduce the alternating moduli paradigm which is argued by its authors in~\cite[Table~2]{TCC:BIPSW18} to provide three orders of magnitude of improvement compared to OPRFs built on top of block ciphers such as \aes{}, {\sf LowMC}~\cite{EC:ARSTZ15}, or {\sf Rasta}~\cite{C:DEGLLL18}.

\subsubsection{From TFHE}

Fully-homomorphic encryption (FHE)~\cite{STOC:Gentry09} is a natural high-end outer component for building round-optimal OPRFs: using FHE, one user would encrypt $x$, let the other user compute an encryption of $\Eval(\msk,x)$ homomorphically, and decrypt the result. Unfortunately, FHE is typically considerably slower than general secure computation, and a significant effort is being invested both in optimizing FHE and in optimizing primitives for homomorphic evaluation via FHE (in the setting of ciphers, standard proposals include {\sf LowMC}~\cite{EC:ARSTZ15}, {\sf FLIP}~\cite{EC:MJSC16}, {\sf Rasta}~\cite{C:DEGLLL18}, {\sf Kreyvium}~\cite{FSE:CCFLNP16}, or {\sf Transistor}~\cite{add:Transistor}, among others). Among the best-performing FHE schemes is {\sf TFHE}~\cite{JC:CGGI20}. It features more compact ciphertexts and relatively fast linear operations, as well as a fast \emph{functional bootstrapping}~\cite{EPRINT:MicPol20,EPRINT:Joye21b} that takes only a few milliseconds.
A simplified circuit model for {\sf TFHE} is the following (also of the soft restriction type, as {\sf TFHE} can in principle evaluate any function).

\begin{description}
\item[Underlying Algebra.] The circuit operates over a small integer ring $\Z_n$ (e.g., $n=2$, $n=17$).
\item[Noise.] A quantity called ``noise'' is increased by every operation, and will lead to an incorrect decryption should it become too high. It is then crucial to prevent it from increasing too much.
\item[Programmable Bootstrap (PBS).] Bootstrapping is a ``refreshing'' operation that brings the noise back down to a given base level. In the case of {\sf TFHE}, bootstrapping is ``programmable'', meaning that it is possible to evaluate an arbitrary function $f:\Z_n\to\Z_n$ implemented by a lookup table without increasing the cost of the bootstrap. This operation is by far the most expensive, so its total number must be minimized---and it is better if those that are used can be evaluated in parallel.
\item[Free Linear Operations.] The evaluation of linear gates of arbitrary indegree is virtually ``free'' in terms of time, but it incurs a slight increase in the noise.
\end{description}

Of course, a linear function over any integer ring cannot be a wPRF, hence any wPRF in the above circuit model must contain at least one functional gate. Therefore, a natural question to answer is the following.

\qbox{Given a small integer $n$ (e.g., $n\leq 17$), are there efficient wPRF candidates computable by arithmetic circuits over $\Z_n$ with additions, multiplications by a constant, and a single layer of indegree-1 programmable bootstrapping gates? What is the smallest number of PBS gates required for a secure wPRF to exist in this model?}

Recently, the work of~\cite{EC:ADDG24} identified that a wPRF of~\cite{TCC:BIPSW18}, while not originally designed with this circuit model in mind, is especially well-suited for evaluation within {\sf TFHE}, as each wPRF evaluation can be computed using solely linear operations together with a \emph{single} layer functional bootstrapping. In fact, it is not too hard to observe that the alternative design of~\cite{TCC:BIPSW18}, denoted \emph{alternative mod-2/mod-3 wPRF} in this writeup, is a (single-bit output) wPRF that uses a \emph{single} $f$-functional gate (we cover this wPRF in Section~\ref{sec:bipsw_original}). This illustrates a general behavior, also observed in other contexts, that shallow wPRFs that optimize to the extreme in some circuit model tend to also perform well in different circuit models (in the context of block ciphers, a similar observation was made for {\sf Kreyvium}~\cite{FSE:CCFLNP16} that performs well for {\sf TFHE} in spite of being originally optimized for BGV-style FHE).


\subsection{Pseudorandom Correlation Generators and Functions}

Pseudorandom correlation generators (PCG) and pseudorandom correlation functions (PCF) are related cryptographic primitives introduced and studied in a recent line of work~\cite{CCS:BCGI18,C:BCGIKS19,CCS:BCGIKR19,C:BCGIKS20,FOCS:BCGIKS20,C:CouRinRag21,C:BCGIKR22,C:BCCD23}. They play an important role in modern secure computation protocols: secure computation considers a scenario where $n$ parties, with respective secret inputs $(x_1, \ldots, x_n)$, wish to jointly compute $f(x_1,\ldots,x_n)$ for some public function $f$ without revealing anything more than the result of the computation. Modern secure computation protocols~\cite{STOC:GolMicWig87,AC:FKOS15,CCS:KelOrsSch16,C:DPSZ12,EC:KelPasRot18} rely on \emph{correlated randomness} to achieve high concrete efficiency: random strings are sampled according to a joint distribution (``correlated'') and are securely distributed among the parties prior to the protocol. Securely generating these strings has long been a major bottleneck in these protocols. PCGs and PCFs enable multiple participants, given short correlated keys, to generate long correlated pseudorandom strings \emph{without any communication}, effectively pushing the bulk of the correlated randomness distribution to an offline computation. Below, we will describe two recent paradigms for constructing PCGs and PCFs from weak pseudorandom functions. For both paradigms, and unlike Section~\ref{subsec:oprf}, the metrics are of the ``hard restriction'' type: the class of circuits handled by the frameworks are highly specific types of very low-depth circuits that can only handle a limited set of gates.


\subsubsection{From Function Secret Sharing}
\label{subsubsec:fss}

Most constructions of PCGs and PCFs (in fact, all those cited above) follow a common template that combines a cryptographic primitive called \emph{function secret sharing} (FSS) with either a pseudorandom generator (for PCGs) or a wPRF (for PCFs). Roughly, an FSS for a function class $\setF$ is a high-level cryptographic primitive that allows sharing any function $f\in \setF$ into a pair of shares $(f_0, f_1)$ such that two parties with a common input $x$ and respective shares $f_0,f_1$ can compute $y_b = \Eval(f_b,x)$ for $b=0,1$ such that $y_0+y_1 = f(x)$. In other words, an FSS allows sharing a function $f\in\setF$ such that from the shares of $f$, two parties can obtain additive shares of $f(x)$ for any input $x$. At the heart of PCG and PCF constructions are variants of the following (semi-formal) theorem.
\begin{theorem}[Theorem~5.3 in~\cite{FOCS:BCGIKS20}]
\label{thm:pcf-ot}
Let $\ring = \{\ring_n\}_{n\in \N}$ denote a family of commutative rings. Let $\PRF=(\KeyGen,\Eval)$ denote a weak pseudorandom function and let $\setF$ denote the family of all functions $x\mapsto c\cdot\Eval(\msk,x)$ for some constant $c$ and PRF key $\msk$. If there exists an FSS for $\setF$, then there exists a PCF for the OT correlation.
\end{theorem}


The ``OT correlation'' (short for oblivious transfer correlation) is, without going into details, the most standard correlation required by secure computation protocols. Hence, the theorem above says that PCFs for useful correlations can be constructed, provided that there are wPRFs that fit into a class of functions for which there exists an FSS scheme.\footnote{More formally, we need \emph{scalar multiples} of wPRFs to fit in the class, that is, functions of the form $x\to c\cdot \Eval(\msk,x)$ for some scalar $c$.} In turn, the class of functions for which we know \emph{efficient} FSS schemes\footnote{Theoretically, there exist advanced constructions of FSS for all polynomial-time functions from strong forms of fully-homomorphic encryption~\cite{C:DHRW16} that are known to exist from the LWE assumption. However, these results are purely of theoretical interest.} is fairly limited~\cite{CCS:BoyGilIsh16}, and corresponds to the set of functions that are linear combinations of ``interval functions'', as formally defined below.

\begin{definition}
Fix a ring $\ring$ and an integer $N \in \mathbb{N}$. Let $\setI = \setI_N(\ring) \subset \ring^{N}$ denote the set of \emph{interval functions} $f_{a,b,\Delta} : [N] \to \ring$, parameterized by $a,b \in [N]$ with $a \leq b$, and $\Delta \in \ring$, and defined as:
\[
f_{a,b,\Delta}(x) = 
\begin{cases}
\Delta & \text{if } x \in [a,b], \\
0 & \text{otherwise}.
\end{cases}
\]

For a given $t \in \mathbb{N}$, define $\setF_{t,N}(\ring)$ as the set of functions of the form
\[
(x_1, \dots, x_t) \mapsto L(f_1(x_1), \dots, f_t(x_t)),
\]
where each $f_i \in \setI$ and $L : \ring^t \to \ring$ is a linear function.
\end{definition}


Equivalently, the class $\setF_{t,N}(\ring)$ can be viewed as the class of functions computable by a depth-$2$ circuit with indegree-1 comparison gates on the bottom (that output a fixed payload if the input is above, or below, a given threshold) and a $t$-ary linear combination gate on top. In the following, we will slightly abuse the definition and view $\setF_{t,N}(\ring)$ both as a class of functions and as the corresponding class of circuits. The following informal lemma summarizes what is known about the existence of efficient FSS schemes for the class $\setF$ (a more formal statement can be found in~\cite{CCS:BoyGilIsh16}).
\begin{lemma}[Informal]
\label{lemma:efficient-fss}
Assume the existence of a length-doubling pseudorandom generator. There exists an FSS for the class $\setF_{t,N}(\ring)$ where the shares $(f_0,f_1)$ of a function $f\in\setF_{t,N}(\ring)$ have size $t\cdot (\secpar\cdot \log N + \log|\ring|)$ and the cost of evaluating a share on an input $x$ is dominated by $2t(\log N + (\log|\ring|)/\secpar)$ invocations of the PRG.
\end{lemma}

In practice, the pseudorandom generator is typically instantiated using {\sf AES}, and the cost of evaluating the FSS becomes dominated by $2t(\log N + (\log|\ring|)/\secpar)$ calls to {\sf AES}, which is extremly efficient (for moderate values of $t$) using modern hardware with support for {\sf AES} instructions. 
FSS constructions for more complex classes are known, but they make a heavy use of public-key cryptography and are considerably less efficient. Therefore, past works on PCGs and PCFs have focussed on developping PRGs and wPRFs compatible with the FSS schemes promised by Lemma~\ref{lemma:efficient-fss}---that is, computable by a depth-$2$ circuit in the class $\setF$ over a suitable ring (typically, the field $\F_{2^\secpar}$, where the reader can think of $\secpar$ as being $128$):

\qbox{Are there weak pseudorandom functions computable by a circuit in $\setF_{t,N}(\F_{2^\secpar})$ for $\lambda=128$ and for some value $t \ll N$?}

In the PCF-from-FSS compiler, $N$ translates to an upperbound on the target number of samples from the correlation (for a wPRF, $N$ must be superpolynomial), and $t$ is a parameter that influences both the PCF key size and the evaluation time. The works of~\cite{FOCS:BCGIKS20,PKC:CouDuc23,C:BCGIKR22} introduced respectively the {\sf VDLPN-wPRF}, 
and the {\sf EALPN-wPRF}, as candidate solutions to the above challenge (with a value $t$ polylogarithmic in $N$). We cover these candidates in Section~\ref{sec:vdlpn} and Section~\ref{sec:ealpn} respectively.

\subsubsection{From Constrained Pseudorandom Functions}

Recently, an alternative paradigm for building PCFs has emerged in~\cite{EC:BCMPR24,AC:CDDKS24}. This paradigm replaces the outer FSS component with a constrained pseudorandom function (CPRF): a pseudorandom function equipped with a special ``constraining'' algorithm that, on the input of a constraint $C$ (a function with binary output from a target class of constraints $\class$) and a PRF key $\msk$, outputs a \emph{constrained key} $\key_C$ that allows evaluating the PRF on all inputs $x$ such that $C(x) = 1$, but leaks no information about the PRF output when $C(x)=0$. The PCFs built from this paradigm enjoy a number of appealing procedures compared to the traditional paradigm: most notably, they come with a \emph{public key setup} (a one-round protocol for securely generating the PCF keys), while efficient (one-round or multi-round) key distribution protocols for earlier PCF constructions are still an open problem. At the heart of this alternative paradigm is the following theorem:

\begin{theorem}
Let $\PRF = (\KeyGen,\Eval)$ denote a weak pseudorandom function with binary output and let $\setF$ denote the family of all functions $x\to \Eval(msk,x)$ and $x\to 1-\Eval(\msk,x)$ for some PRF key $\msk$. If there exists a CPRF for the class of constraints $\setF$, then there exists a PCF for the OT correlation.
\end{theorem}

The work of~\cite{EC:BCMPR24} additionally introduced an efficient CPRF construction where the cost of evaluating the CPRF (and therefore the PCF) is dominated by a single exponentiation over an elliptic curve (on a modern laptop and for the fastest {\sf OpenSSL} curves, this can be done in about $8\mu$s). This CPRF builds upon the seminal Naor-Reingold PRF~\cite{FOCS:NaoRei97}, and is restricted to constraints computable by an \emph{inner-product membership} (IPM) predicate.
\begin{definition}
Fix a security parameter $\secpar$, integers $B,n \in \N$, and let $g:\bit^\secpar \to [B]^n$ be a function and $S \subset \N$ be a finite set. Let $\IPM = \IPM_{B,n}(S,g)$ denote the set of \emph{inner-product membership predicates over $S$}, i.e., the functions $P_{y}:[B]^n\to \bit$ for $y \in [B]^n$ defined as follows:
\[P_{y}:x\mapsto \begin{cases} 1 \text{ if } \langle g(x),y\rangle \in S \algcomment{The inner product is computed over $\N$}\\
0 \text{ otherwise}~.
\end{cases}\]
\end{definition}
In other words, after fixing some public ``preprocessing'' function $g$ on the input, and a subset $S$, an IPM predicate is a depth-$3$ circuit with a $g$-gate at the bottom (that takes a single input $x$ and has $n$ outputs), an arbitrary $n$-ary linear combination gate in the middle, and an indegree-1 membership-test gate on top. A formal statement of the following informal lemma can be found in~\cite{EC:BCMPR24}.
\begin{lemma}[Informal]
\label{lemma:efficient-cprf}
Assume the power-DDH assumption over a group $\G$ of prime order $p$. Then there exists a CPRF for the class of constraints $\IPM = \IPM_{B,n}(S,g)$ where the size of a constrained key is $(n+|S|)\cdot \log|\G|$, and the cost of evaluating the CPRF is dominated by $n$ multiplications over $\Z_p$ and one exponentiation over $\G$.
\end{lemma}

Above, the number of multiplications can be reduced to the Hamming weight of $g(x)$. Then, to obtain efficient PCFs, it remains to find wPRFs that can be computed by an inner-product membership predicate, called IPM-wPRF in~\cite{EC:BCMPR24}:

\qbox{Are there weak pseudorandom functions computable by a circuit from $\IPM_{B,n}(S,g)$ where $B$ is some integer bound and $S$ is a small subset of integers?}

In the PCF-from-CPRF compiler, a larger $S$ translates to a larger key size (as it grows linearly with $|S|$). The work of~\cite{EC:BCMPR24} showed that two wPRF from previous works are actually IPM-wPRF: the Boneh-Ishai-Passelègue-Sahai-Wu , or alternating moduli wPRF~\cite{TCC:BIPSW18} and the Goldreich-Applebaum-Raykov wPRF~\cite{DBLP:journals/eccc/ECCC-TR00-090,TCC:AppRay16b}. We cover both candidates in Section~\ref{sec:bipsw} and Section~\ref{sec:goldreich} respectively.

\subsection{VOLE-Based Zero-Knowledge Proofs}

Zero-knowledge proofs (ZKP) are fundamental cryptographic primitives that allow proving properties about a system that depends on secret values without revealing these values. While they were mainly an object of theoretical interest after their introduction in~\cite{GolMicRac89}, four decades of research have led to a plethora of highly practical ZKP that are now routinely used in real-world systems, from authentication to anonymous blockchains and cryptocurrencies (a recent example is Google's wallet~\cite{googlewallet}).

\paragraph{Definition.} An $\NP$-language is a set of strings $\Lang\subset \bit^*$ such that there is an algorithm $\Rel_\Lang$ (the \emph{relation}) that can verify proofs of membership to $\Lang$ in polynomial-time: a string $x$ is in $\Lang$ if and only if there exists a \emph{witness} $w$, of size polynomial in $x$, such that $\Rel_\Lang(x,w)$ returns ``accept'' in time polynomial in $|x|$. A zero-knowledge proof for an $\NP$-language $\Lang$ is a two-party secure computation protocol between a prover and a verifier with the following characteristics:
\begin{itemize}
	\item Both parties have as common input a word $x \in \Lang$, and the (honest) prover additionally holds a \emph{witness} for $x$---that is, a string $w$ such that $\Rel_\Lang(x,w)=1$.
	\item At the end of the interaction, the verifier outputs 1 (``accept'') or 0 (``reject'').
	\item An honest prover (with a witness $w$) should always cause the verifier to accept. However, if $x\notin\Lang$ (hence no valid witness $w$ exists), no prover can possibly cause the verifier to output 1 with non-negligible probability.
	\item Eventually, the proof is \emph{zero-knowledge}, meaning that it leaks no information about $w$ (beyond the fact that $x\in \Lang$). This is formalized by requiring that the view of the verifier (the transcript of its interaction with the prover) could have been simulated without the witness $w$, and the simulation is indistinguishable from the real transcript (hence, in essence, the verifier cannot tell the transcript apart from a transcript computed independently of $w$).
\end{itemize}

\paragraph{Efficient zero-knowledge proof systems.} Modern zero-knowledge proof systems require tradeoffs and compromises. Typically, they will handle restricted types of algebraic statements~\cite{C:Schnorr89,EC:GroSah08}, or be fully generic and optimal on almost any aspect, but incur an often prohibitive computational overhead on the prover side~\cite{EC:BCCGP16,EC:Groth16,CCS:AHIV17,SP:BBBPWM18,ICALP:BBHR18,EC:BCRSVW19,CCS:MBKM19,C:Setty20,EC:CBBZ23} (the citations are a short sample from the vast SNARK/STARK ecosystem of proofs). Building on the seminal work of~\cite{STOC:IKOS07}, a recent line of work has put forth \emph{VOLE\footnote{Short for Vector Oblivious Linear Evaluation: a mechanism that distributes vectors $\vec u, \vec v$ to the prover and $\Delta, \vec u+\Delta\cdot \vec v$ to the verifier, where $\Delta$ acts as a MAC to ``authenticate'' the vector $\vec v$.}-based zero-knowledge} as a promising middle ground~\cite{CCS:BCGI18,SP:WYKW21,ITC:DitIshOst21,CCS:YSWW21,C:BMRS21}: VOLE-based ZKP can handle arbitrary statements, have reasonable proof sizes (though higher than SNARKs/STARKs for large statements), and small prover costs, confering them an overall highly competitive scalability. VOLE-based ZKP are deployed in real-world applications (e.g.\ TLSNotary~\cite{tlsnotary}), and in commercial solutions (such as {\sf zkPass}~\cite{zkpass} or {\sf zkTLS}~\cite{primus}).

A typical ``pain point'' of zero-knowledge proofs are \emph{mixed statements}: proofs of statements that mix arithmetics over different structures (a typical example is neural network training, that often mixes arithmetic over large fields with Boolean operations). Mixed statements have been the focus of a significant attention, both in general secure computation~\cite{NDSS:DemSchZoh15,SP:GYKW24,EC:BCGGIK21} and in zero-knowledge proofs~\cite{C:ChaGanMoh16,CCS:BBMRS21,EPRINT:OKMZ24,DBLP:conf/eurocrypt/AgarwalBBS25}. A recent work~\cite{DBLP:conf/eurocrypt/AgarwalBBS25} observed that shallow wPRFs provide an ideal tool to handle mixed arithmetic in VOLE-based zero-knowledge proofs. At the high level, the key idea is the following: in VOLE-based zero-knowledge, the secret witness is authenticated to the verifier via an information-theoretic homomorphic MAC. Previous works had observed that when the secret witness must be used over two different structures, e.g., the field $\F_2$ and a larger field $\F_p$, it suffices to have pairs of identical bits MAC-authenticated both in $\F_2$ and $\F_p$ to securely ``switch'' between the two structures. The work of~\cite{DBLP:conf/eurocrypt/AgarwalBBS25} observes that an efficient way to achieve this is by first authenticating a short wPRF key over both fields and proving the equality. This is a mixed statement, so it can be expensive, but it is crucially independent of the size of the full statement since the key is short: it will get amortized away. Then, identical \emph{pseudorandom} bits are generated over both fields by evaluating the wPRF and proving (separately over each field) the correct evaluation. To make this procedure efficient, the key is to find a wPRF that admits short proofs of correct evaluation \emph{both over $\F_2$ and $\F_p$}.

\paragraph{Cost model.} The basic cost model of modern VOLE-based zero-knowledge proofs is the usual ``pay-per-product'' MPC cost model, inherited from {\sf Quicksilver}~\cite{CCS:YSWW21} (the underlying proof system used in these works): when proving that an arithmetic circuit $C$ (with additions and multiplications over a field $\F$) evaluates to a certain value on a secret witness $w$, all linear operations (additions, multiplications by a constant) are for free (they do not incur any overhead in the proof size), while multiplications increase the proof size. Therefore, the wPRF must be of the standard ``MPC-friendly'' type: its circuit description must minimize the number of multiplications.

In the setting of mixed statements, however, a subtlety arises: given a statement mixing operations over two fields $(\F,\F')$, the same wPRF must be computed by an arithmetic circuit over each of these fields. One must therefore be careful, as a low multiplication-gate count over one field does not typically translate to a low multiplication-gate count over the other field---in fact, the opposite is often true: linear operations over a field $\F_p$ are highly nonlinear over a field $\F_q$ when $p,q$ are coprime~\cite{Raz87,Smo87} (this behavior is actually a crucial component of the alternating moduli paradigm, that we cover in Section~\ref{sec:bipsw}). However, when $\F,\F'$ are prime order fields, a simple observation is that it suffices to have a wPRF computable by a circuit with a small number of multiplications \emph{over the integers}\footnote{The circuit operates directly on the input bits of the wPRF, and must output integers that are guaranteed to be bits as well.}: the natural arithmetization of this circuit (replacing additions and multiplications over the integers with the field operations) yields a valid arithmetic circuit for the wPRF over any prime-order field.
 In turn, any wPRF computable by a $d$-\emph{local} Boolean circuit (i.e., such that each output bit depends on at most $d$ input bits) can be arithmetized into an arithmetic circuit over the integers with at most $d$ multiplications. Combining these two observations, we get the following target cost model.
 \begin{description}
\item[Underlying Algebra.] The circuit operates over the Boolean field $\F_2$.
\item[Locality.] The circuit is ``$d$-local'' for some small quantity $d$: every output bit depends on at most $d$ input bits.
\end{description}

The question below summarizes the quest for the best possible candidate:

\qbox{Given a target bound $n(\secpar)$ on the key size (for a security parameter $\secpar$) and a target number $N(\secpar)$ of samples, what is the smallest $d$ such that there exist weak pseudorandom functions with key size $n(\secpar)$ computable by a $d$-local Boolean circuit?}

The study of local PRGs and wPRFs has a rich history in cryptography. We cover the main candidates in Section~\ref{sec:goldreich}.



